\section*{Приложение 1. Модульная спецификация и алгоритмическое описание}
\addcontentsline{toc}{section}{Приложение 1. Модульная спецификация и алгоритмическое описание}
\renewcommand{\thetable}{П1.\arabic{table}}
\setcounter{table}{0}

\subsection*{П1.1. Модульная спецификация пакета \texttt{app}}

Пакет \texttt{app} (модуль \texttt{uir-go-profiler/\allowbreak app}) состоит из одного
файла \texttt{processor.go}, определяющего один тип данных и шесть
экспортируемых функций обработки. Полная спецификация — таблица~\ref{tab:p1.1}.

\begin{table}[htbp]
\caption{Модульная спецификация пакета \texttt{app}}
\label{tab:p1.1}
\centering
\begin{tabular}{p{3.1cm}p{6.6cm}p{4.3cm}}
\toprule
Имя & Сигнатура & Назначение \\
\midrule
\texttt{Record} & \texttt{struct\{ID int; Name string; Category string; Amount float64; Active bool\}} & доменная сущность, обрабатываемая всеми функциями пакета \\
\texttt{ParseRecords} & \texttt{func([]byte) ([]Record, error)} & декодирует JSON-массив записей; возвращает обёрнутую ошибку при некорректном JSON \\
\texttt{FilterActive} & \texttt{func([]Record) []Record} & возвращает подмножество записей с \texttt{Active == true}, сохраняя порядок \\
\texttt{FindByID} & \texttt{func([]Record, int) (Record, bool)} & линейный поиск записи по \texttt{ID}; второе значение — признак найдена/нет \\
\texttt{Aggregate} & \texttt{func([]Record) map[string]float64} & суммирует \texttt{Amount} по значению \texttt{Category} \\
\texttt{FormatNames} & \texttt{func([]Record) []string} & форматирует каждую запись в строку вида \texttt{"Имя (Категория): \$Сумма"} \\
\texttt{Deduplicate} & \texttt{func([]Record) []Record} & удаляет записи с повторяющимся \texttt{ID}, сохраняя первое вхождение \\
\bottomrule
\end{tabular}
\end{table}

Все функции — чистые (без побочных эффектов, без обращения к глобальному
изменяемому состоянию, за единственным намеренным исключением в
диагностическом сценарии \texttt{regression/\allowbreak memory-growth}, где вводится
пакетная переменная \texttt{aggregateHistory} для демонстрации утечки
памяти, раздел~\ref{sec:5.1}), что обеспечивает воспроизводимость измерений
между запусками бенчмарков.

\subsection*{П1.2. Алгоритмическое описание критерия обнаружения регрессии}

Представление выбрано единообразно для всех алгоритмов данного приложения —
псевдокод в нотации \texttt{algorithmicx} (императивный стиль, близкий к
Go). Алгоритм~\ref{alg:criterion} — псевдокод критерия~\eqref{eq:criterion}
из раздела~\ref{sec:2.2}, реализованного в \texttt{scripts/\allowbreak compare.sh} и
\texttt{scripts/\allowbreak lib.sh}.

\begin{algorithm}[htbp]
\caption{Обнаружение регрессии производительности}
\label{alg:criterion}
\begin{algorithmic}[1]
\Function{DetectRegression}{$baseline$, $pkg$, $T$, $\alpha$, $count$}
  \State $current \gets \Call{RunBenchmarks}{pkg, count}$ \Comment{go test -bench -count=$count$}
  \State $report \gets \Call{Benchstat}{baseline, current, \alpha}$ \Comment{CSV: $\Delta_{b,m}$, значимость по $\alpha$}
  \State $regressed \gets \textbf{false}$
  \ForAll{$(b, m, \Delta_{b,m}) \in report$} \Comment{только значимые на уровне $\alpha$ строки}
    \If{$\Delta_{b,m} > T$}
      \State $regressed \gets \textbf{true}$
      \State \Call{AppendReportRow}{$b$, $m$, $\Delta_{b,m}$}
    \EndIf
  \EndFor
  \If{$regressed$}
    \State \Return FAIL \Comment{код завершения 1 --- блокирует merge}
  \Else
    \State \Return OK
  \EndIf
\EndFunction
\end{algorithmic}
\end{algorithm}

Алгоритм~\ref{alg:dedup} сопоставляет две реализации \texttt{Deduplicate}
(эталонную, $O(n)$, и внесённую сценарием \texttt{regression/\allowbreak quadratic-dedup},
$O(n^2)$, раздел~\ref{sec:2.2}) в едином представлении, чтобы разница в
сложности была видна непосредственно из структуры псевдокода (вложенный
цикл по уже обработанному префиксу вместо обращения к хеш-таблице).

\begin{algorithm}[htbp]
\caption{Deduplicate --- эталонная ($O(n)$) и регрессионная ($O(n^2)$) реализации}
\label{alg:dedup}
\begin{algorithmic}[1]
\Function{Deduplicate\_Baseline}{$records$} \Comment{$O(n)$}
  \State $seen \gets \emptyset$ \Comment{хеш-множество, амортизированный $O(1)$ доступ}
  \State $result \gets []$
  \For{$r \in records$}
    \If{$r.ID \notin seen$}
      \State $seen \gets seen \cup \{r.ID\}$
      \State $result.\Call{Append}{r}$
    \EndIf
  \EndFor
  \State \Return $result$
\EndFunction
\Statex
\Function{Deduplicate\_Quadratic}{$records$} \Comment{$O(n^2)$ --- сценарий regression/quadratic-dedup}
  \State $result \gets []$
  \For{$i \gets 0 \dots n-1$}
    \State $duplicate \gets \textbf{false}$
    \For{$j \gets 0 \dots i-1$} \Comment{линейный поиск по уже обработанному префиксу}
      \If{$records[j].ID = records[i].ID$}
        \State $duplicate \gets \textbf{true}$; \textbf{break}
      \EndIf
    \EndFor
    \If{\textbf{not} $duplicate$}
      \State $result.\Call{Append}{records[i]}$
    \EndIf
  \EndFor
  \State \Return $result$
\EndFunction
\end{algorithmic}
\end{algorithm}

\clearpage
\section*{Приложение 2. Исходный код прототипа}
\label{app:code}
\addcontentsline{toc}{section}{Приложение 2. Исходный код прототипа}

Полный исходный код доступен в репозитории прототипа; ниже приведены все
файлы, определяющие поведение системы (тестовое приложение, слой сравнения,
CI-пайплайн). Все функции пакета \texttt{app} снабжены комментариями в
формате godoc непосредственно перед объявлением (стандарт комментирования
экосистемы Go, аналог Javadoc/Doxygen); скрипты слоя сравнения
комментированы построчно перед каждым содержательным блоком.

\subsection*{П2.1. \texttt{app/\allowbreak processor.go} --- бизнес-логика тестового приложения}
\lstinputlisting[language=Go]{../app/processor.go}

\subsection*{П2.2. \texttt{app/\allowbreak processor\_test.go} --- unit-тесты}
\lstinputlisting[language=Go]{../app/processor_test.go}

\subsection*{П2.3. \texttt{app/\allowbreak bench\_test.go} --- бенчмарки}
\lstinputlisting[language=Go]{../app/bench_test.go}

\subsection*{П2.4. \texttt{scripts/\allowbreak lib.sh} --- реализация критерия регрессии}
\lstinputlisting[language=bash]{../scripts/lib.sh}

\subsection*{П2.5. \texttt{scripts/\allowbreak compare.sh} --- решающий шаг}
\lstinputlisting[language=bash]{../scripts/compare.sh}

\subsection*{П2.6. \texttt{scripts/\allowbreak profile-diff.sh} --- диагностический шаг (pprof)}
\lstinputlisting[language=bash]{../scripts/profile-diff.sh}

\subsection*{П2.7. \texttt{scripts/\allowbreak capture-baseline-profiles.sh} --- генерация эталонных профилей}
\lstinputlisting[language=bash]{../scripts/capture-baseline-profiles.sh}

\subsection*{П2.8. \texttt{scripts/\allowbreak experiments/\allowbreak sensitivity.sh} --- эксперимент раздела 5.2 (чувствительность порога)}
\lstinputlisting[language=bash]{../scripts/experiments/sensitivity.sh}

\subsection*{П2.9. \texttt{scripts/\allowbreak experiments/\allowbreak power\_analysis.sh} --- эксперимент раздела 5.2 (мощность критерия)}
\lstinputlisting[language=bash]{../scripts/experiments/power_analysis.sh}

\subsection*{П2.10. \texttt{scripts/\allowbreak experiments/\allowbreak overhead.sh} --- эксперимент раздела 5.3 (накладные расходы)}
\lstinputlisting[language=bash]{../scripts/experiments/overhead.sh}

\subsection*{П2.11. \texttt{.github/\allowbreak workflows/\allowbreak perf-check.yml} --- CI-пайплайн}
\lstinputlisting[basicstyle=\ttfamily\fontsize{9}{10.5}\selectfont]{../.github/workflows/perf-check.yml}
